\chapter{Implementation}

\section{Module-wise Implementation}
Provide a detailed breakdown of implementation for each module:
\begin{itemize}
\item Module name and its purpose
\item Key functionalities implemented
\item Data structures used
\item External dependencies or libraries
\item Implementation challenges and solutions
\end{itemize}

\section{Code Snippets and Explanation}
Include important code snippets with explanations:
\begin{itemize}
\item Key algorithms with comments
\item Critical functions or methods
\item Initialization and configuration code
\item Error handling mechanisms
\end{itemize}

Example code structure:
\begin{verbatim}
# Example: Authentication Module
def authenticate_user(username, password):
    """
    Authenticates user credentials.
    Args:
        username: User's login identifier
        password: User's login password
    Returns:
        Boolean indicating authentication success
    """
    # Validate input parameters
    if not username or not password:
        raise ValueError("Invalid credentials")
    
    # Database lookup and password verification
    user = db.find_user(username)
    if user and verify_password(password, user.hashed_pwd):
        return True
    return False
\end{verbatim}

\section{Database Implementation}

\subsection{Schema and Structure}
Describe the database schema, tables, and their relationships.

\subsection{SQL Table Structures}
Provide CREATE TABLE statements for all database tables with appropriate data types, constraints, and indexes.

\subsection{ER Diagram Reference}
Link to ER diagram showing complete database design.

\begin{table}[htbp]
\centering
\caption{Database Tables Overview}
\begin{tabular}{p{0.2\textwidth}p{0.3\textwidth}p{0.4\textwidth}}
\toprule
Table Name & Primary Key & Key Attributes \\
\midrule
users & user\_id & username, email, password\_hash \\
projects & project\_id & name, description, created\_by \\
tasks & task\_id & title, status, assigned\_to, deadline \\
\bottomrule
\end{tabular}
\end{table}

\clearpage
